\section*{Chapter \ref{ch:rs:performance-measurement} --- Performance Measurement}

\begin{solution}{exo:rs:performance-measurement:1}
About $\sqrt{252/1\,260} = 0.45$ (the $\mathrm{SR}^2/2$ term is negligible per day); the 95\% interval runs from 0.12 to 1.88.
\end{solution}

\begin{solution}{exo:rs:performance-measurement:2}
$0.144/0.336 = 0.43$. The \omterm{def:rs:performance-measurement:drawdown}{maximum drawdown} is one draw from a wide distribution: the same strategy's five-year maximum exceeds 30\% in 14\% of paths and stays
far below it in most; a ratio with one extreme in its denominator inherits that instability.
\end{solution}

\begin{solution}{exo:rs:performance-measurement:3}
Over $n$ periods, $\sum r^+ / \sum r^- = (\frac1n\sum r^+)/(\frac1n\sum r^-) = \mathbb{E}_n[(r - 0)^+]/\mathbb{E}_n[(0 - r)^+]$, the sample \omterm{def:rs:performance-measurement:sortino}{Omega ratio} at zero.
\end{solution}

\begin{solution}{exo:rs:performance-measurement:4}
$\sum\theta_j^2 = 0.25 + 0.09 + 0.04 = 0.38$; volatility factor $\sqrt{0.38} = 0.616$; reported Sharpe ratio $0.6/0.616 = 0.97$;
$\rho_1 = (0.5 \times 0.3 + 0.3 \times 0.2)/0.38 = 0.55$, $\rho_2 = 0.5 \times 0.2/0.38 = 0.26$.
\end{solution}

\begin{solution}{exo:rs:performance-measurement:5}
$12/\sqrt{12 + 2(11 \times 0.55 + 10 \times 0.26)} = 2.22$ against $\sqrt{12} = 3.46$: the naive \omterm{def:rs:performance-measurement:annualisation}{annualisation} overstates the Sharpe ratio by a factor of
$3.46/2.22 = 1.56$, 56\%.
\end{solution}

\begin{solution}{exo:rs:performance-measurement:6}
The beta averages a kinked payoff: over 100 simulated years the daily beta is 0.98 on the index's down days and 0.71 on up days, and in a crash the put's delta
rises towards one, so the loss is far larger than 0.45 times the index's fall. The displayed history (0.59 and 0.54) holds no crash to reveal the kink.
\end{solution}

\begin{solution}{exo:rs:performance-measurement:7}
Within one year: 0.6\% for the trend strategy, 11\% for short volatility; within ten years: 31\% and 65\%. The short-volatility strategy's losses come in jumps
(a crash month takes a large share of capital at once), so a 30\% fall does not need the slow accumulation of bad days a diffusion needs.
\end{solution}

\begin{solution}{exo:rs:performance-measurement:8}
The fund's \omterm{def:rs:performance-measurement:sortino}{Sortino ratio} (1.87 in the chapter's example) and its record come from a history without a crash; a downside measure cannot see a downside the
sample lacks. Its model gives a 41\% chance of a 30\% drawdown in five years against the trend fund's 14\%, and over 100 years its skewness is $-1.23$.
Measure the strategy's exposure (a short put) and stress it, not the history.
\end{solution}

\begin{solution}{pb:rs:performance-measurement:1}
\begin{enumerate}
\item 1.00 and 1.09, each with an iid standard error of 0.45 (HAC 0.46 and 0.45).
\item The standard error depends on the length of the history in years, about $\sqrt{q/T}$, and all four span five years.
\item 0.63 against 0.46: serially correlated returns carry less information than their number suggests.
\item 4.04 with a standard error of 0.45 (HAC 0.52): a real edge, but nothing about capacity, crowding or what an adverse day can cost at size.
\item 7.4\% reported, 11.2\% true.
\item 0.99 reported, 0.61 with Lo's correction, 0.59 unsmoothed, 0.66 true.
\item $(0.43, 0.36, 0.21)$, from first autocorrelations of 0.64 and 0.25.
\item Its returns are autocorrelated, so the volatility of a year's sum is larger than $\sqrt{12}$ times the monthly volatility.
\item Trend $-33.6\%$ and 791 days; short volatility $-14.8\%$ and 48 days; smoothed $-11.0\%$ and 22 months; market maker $-4.2\%$ and 72 days.
\item The short-volatility history on all three (Calmar 0.78, Sortino 1.87, Omega 1.49 against 0.43, 1.48, 1.17); it also has the worse tail.
\item $+5.38$ and 102.5 in the history; $-1.23$ and 121 over 100 years.
\item 9\%.
\item 48 trading days; a one-way turnover of 1.86\% of capital a day.
\item \textbf{Named result.} A 30\% drawdown within five years: 14\% for the trend strategy, 41\% for short volatility. The smoothed book's Sharpe ratio has a
standard error of 0.63 with its autocorrelation (HAC) against 0.46 assuming independence.
\item Trend: 0.83 and $-37\%$. Short volatility: 0.62 and $-72\%$.
\item 0.45 (standard error 0.09) in the history; down-day and up-day betas of 0.59 and 0.54 in the history, 0.98 and 0.71 over 100 years.
\item The trend strategy, or the short-volatility one only sized to its crash (a stress loss within the drawdown budget), not to its history's volatility.
\item Its standard error (iid and HAC), the length of the history, the number of trials behind it, and the drawdown probability of the strategy's model.
\item Sell options or other tail risk, smooth marks, or choose the reporting window; a manipulation-proof measure or a stress test answers each.
\item That the strategy's true Sharpe ratio is probably between 0.1 and 1.9, and nothing about its tail.
\end{enumerate}
\end{solution}

\begin{solution}{iq:rs:performance-measurement:1}
The standard error is about $\sqrt{252/756} = 0.58$, so the estimate is 1.7 standard errors from zero: fairly but not very confident, before any correction
for the number of strategies tried (chapter~20).
\end{solution}

\begin{solution}{iq:rs:performance-measurement:2}
Sell out-of-the-money options or other tail risk (steady gains, rare large losses), smooth the marks of illiquid positions, or pick the window. Look at the
shape across a long simulated or stressed history, the down-market beta, the autocorrelation of returns, and the positions themselves; Goetzmann and co-authors
give a measure that cannot be gamed.
\end{solution}

\begin{solution}{iq:rs:performance-measurement:3}
It inflates the annualised Sharpe ratio (stale marks lower the measured volatility) and makes it less precise. Correct with Lo's autocorrelation-adjusted
\omterm{def:rs:performance-measurement:annualisation}{annualisation}, unsmooth the returns with the estimated profile of Getmansky, Lo and Makarov, and use HAC standard errors.
\end{solution}

\begin{solution}{iq:rs:performance-measurement:4}
For a Brownian motion with positive drift it grows logarithmically with the horizon (square root with no drift). A \omterm{def:rs:vectorised-backtests:backtest}{backtest}'s maximum is a single draw, and
it cannot contain the events the sample lacks; use the distribution of drawdowns under a model and stress tests.
\end{solution}

\begin{solution}{iq:rs:performance-measurement:5}
The loss distribution of its bad days (the 1.5\% adverse-selection days), capacity and how the edge scales with size, its inventory limits, the correlation of its
losses with market stress, and whether the P\&L decomposes into spread capture and adverse selection that make sense (chapter~23).
\end{solution}

\begin{solution}{iq:rs:performance-measurement:6}
With $\hat\mu$ and $\hat\sigma^2$ independent for normal data, $\operatorname{Var}\hat\mu = \sigma^2/T$ and $\operatorname{Var}\hat\sigma^2 = 2\sigma^4/T$. The delta
method on $\mathrm{SR} = \mu/\sigma$, with derivatives $1/\sigma$ and $-\mu/(2\sigma^3)$, gives $\operatorname{Var}\widehat{\mathrm{SR}} = 1/T + \mu^2/(2\sigma^2 T) =
(1 + \mathrm{SR}^2/2)/T$.
\end{solution}
